\subsection{相似变换群。}

设 \(\Phi\) 是 E 上的一个埃尔米特型。A-模 E 的一个自同构 \(u\) 称为一个相似变换（关于 \(\Phi\) 的），如果存在 A 的一个可逆元素 \(\alpha\) 使得

\[
\Phi (u (x), u (y)) = \alpha \Phi (x, y)\tag{7}
\]

对 E 中所有 x、y 成立。相似变换构成一个群 Γ。

当 \(\Phi\) 取的值都是 A 的正则元素时（例如当 A 是一个体且 \(\Phi \neq 0\) 时），满足 (7) 的 A 的元素 \(\alpha\) 由 \(u\) 唯一确定；称之为相似变换 \(u\) 的乘子。在 (7) 中把 \(x\) 换成 \(\lambda x\) ，我们就看出 \(\alpha\) 属于 A 的中心；在 (7) 中交换 \(x\) 与 \(y\) ，我们进一步看出 \(\overline{\alpha} = \alpha\) 。若对 \(u \in \Gamma\) 以 \(\alpha(u)\) 表示 \(u\) 的乘子，则映射 \(u \to \alpha(u)\) 是从 \(\Gamma\) 到 A 的中心的可逆元素所成的乘法群中的一个同态。这个同态的核是与 \(\Phi\) 相伴的酉群，因而它是 \(\Gamma\) 的一个正规子群。设 \(\beta\) 是 A 的中心的一个可逆元素，\(v\) 是比为 \(\beta\) 的位似，\(w\) 是 E 的一个酉自同构；则 \(vw = wv\) 是 E 的一个相似变换，其乘子是 \(\beta\overline{\beta}\) 。反之，设 \(u\) 是一个相似变换，其乘子具有形式 \(\beta\overline{\beta}\) （ \(\beta\) 表示 A 的中心的一个可逆元素）；则 \(uv^{-1}\) 是一个酉自同构 \(w\) ，从而 \(u\) 具有 \(vw\) 的形式。

现在假设 A 是一个体，E 是有限维向量空间，且 \(\Phi\) 非退化。对每个相似变换 \(u\) （乘子为 \(\alpha\) ），我们有

\[
\Phi (x, \alpha y) = \alpha \Phi (x, y) = \Phi (u (x), u (y)) = \Phi (x, u ^ {*} (u (y))),
\]

因此 \(u^*u\) 是比为 \(\alpha\) 的位似。若 A 是交换的，且以 \(n\) 表示 E 的维数，则由上式及 § 1 第 10 小节的公式 (50) 得出

\[
(\det u) \overline {{(\det u)}} = \alpha^ {n}.\tag{8}
\]

于是我们区分两种情形：

1°) 整数 n 是奇数，即 \(n = 2q + 1\) 。此时令 \(\rho = \alpha^{-q}(\det u)\) ，就有 \(\alpha = (\det u) \overline{(\det u)} \alpha^{-2q} = \rho \bar{\rho}\) 。于是 u 是比为 \(\rho\) 的位似与一个酉自同构之积。

2°) 整数 n 是偶数，设 n = 2q。此时令 \(\rho = \alpha^{-q}(\det u)\) ，就有 \(\rho\bar{\rho} = 1\) 。特别地，当 J 是恒等映射时，有 \((\det u)^{2} = \alpha(u)^{2q}\) ；使 det \(u = \alpha(u)^{q}\) （相应地，det \(u = -\alpha(u)^{q}\) ）的相似变换 u 称为正向的（相应地，逆向的）；正向相似变换构成 \(\Gamma\) 的一个指数为 2 的正规子群；

比 ≠ 0 的位似是正向相似变换；行列式为 1 的正交变换（第 2 小节）也是如此；行列式为 -1 的正交变换是逆向相似变换。

前面的定义与结果对 ε-埃尔米特型（§ 3，第 1 小节）仍然有效，特别对交错型也有效。

设 A 是一个交换域，Q 是 E 上的一个二次型 \(\neq 0\) 。我们称 E 的满足下述条件的自同构 \(u\) 为（关于 Q 的）相似变换：存在 A 的一个非零元素 \(\alpha\) （称为 \(u\) 的乘子），使得有 \(Q(u(x)) = \alpha Q(x)\) （对每个 \(x \in E\) ）。显然此时 \(u\) 是关于与 Q 相伴的双线性型的、乘子为 \(\alpha\) 的相似变换；当 A 的特征 \(\neq 2\) 时逆命题成立。
% semantic-fragments: legacy-input-seam
\relax{}
